% GENERATED FILE. Edit canonical lessons, then run npm run generate:books.
% canonical-source-hash: fnv1a64:ead8ca68cf2bbe9a
% canonical-lessons: ML-C195-ennum, ML-C195-tiyati, ML-C195-varantyam, ML-C195-pettennu, ML-C195-munkutti

\chapter{Every day, Date, Weekend, Quickly, In advance}
\label{ch:ml-every-day-date-weekend-quickly-in-advance}
\hlchaptermodality{195}

\begin{chapteropening}
\textbf{By the end of this chapter:} \emph{I can say every day, a date, the weekend, quickly and in advance.}

Complete the last lesson of chapter 195: I can say every day, a date, the weekend, quickly and in advance.
\end{chapteropening}

\section[\texorpdfstring{\ml{എന്നും} (ennuṁ)}{എന്നും (ennuṁ)}]{\ml{എന്നും} (ennuṁ) — every day, always}

\label{lesson:ML-C195-ennum}



\begin{quote}
Before the new one: say the Malayalam for health, then the Malayalam for an eyelid.
\end{quote}

\subsection*{You'll want to know: \ml{എന്നും}}
\textbf{\ml{എന്നും}} — \emph{ennuṁ} — \textquotedblleft{}every day, always\textquotedblright{}.

\begin{cousinweb}
\textbf{\ml{എന്നും}} (\emph{ennuṁ}) is native Dravidian: DEDR 5151 lists it, \textquotedblleft{}always\textquotedblright{}, under the old question stem, with Tamil \emph{yā}: \textquotedblleft{}on whatever day\textquotedblright{}.
\end{cousinweb}

\begin{grammarlens}[title={What you've built}]
One more word for this chapter.
\end{grammarlens}

\subsection*{Your turn}
\begin{itemize}
\raggedright
  \item \emph{Say it:} \emph{ennuṁ}
  \item \emph{Say it:} \emph{ennuṁ}, once more
  \item \emph{Say it:} say \emph{ennuṁ}
  \item \emph{Recall:} say the Malayalam for health, then the Malayalam for an eyelid, then say \emph{ennuṁ} again
\end{itemize}

\begin{tcolorbox}[breakable,colback=teal!4,colframe=teal!35!black,title={Before you move on}]
What does \ml{എന്നും} mean? (\textquotedblleft{}Every day, always\textquotedblright{}.) Say it once more.
\end{tcolorbox}
\section[\texorpdfstring{\ml{തീയതി} (tīyati)}{തീയതി (tīyati)}]{\ml{തീയതി} (tīyati) — a date}

\label{lesson:ML-C195-tiyati}



\begin{quote}
Before the new one: say the Malayalam for an eyelid, then the Malayalam for every day.
\end{quote}

\subsection*{You'll want to know: \ml{തീയതി}}
\textbf{\ml{തീയതി}} — \emph{tīyati} — \textquotedblleft{}a date\textquotedblright{}.

\begin{cousinweb}
\textbf{\ml{തീയതി}} (\emph{tīyati}) is a tadbhava of Sanskrit \emph{tithi}, \textquotedblleft{}a lunar day\textquotedblright{} (CDIAL 5810), reshaped in the mouth.
\end{cousinweb}

\begin{grammarlens}[title={What you've built}]
One more word for this chapter.
\end{grammarlens}

\subsection*{Your turn}
\begin{itemize}
\raggedright
  \item \emph{Say it:} \emph{tīyati}
  \item \emph{Say it:} \emph{tīyati}, once more
  \item \emph{Say it:} say \emph{tīyati}
  \item \emph{Recall:} say the Malayalam for an eyelid, then the Malayalam for every day, then say \emph{tīyati} again
\end{itemize}

\begin{tcolorbox}[breakable,colback=teal!4,colframe=teal!35!black,title={Before you move on}]
What does \ml{തീയതി} mean? (\textquotedblleft{}A date\textquotedblright{}.) Say it once more.
\end{tcolorbox}
\section[\texorpdfstring{\ml{വാരാന്ത്യം} (vārāntyaṁ)}{വാരാന്ത്യം (vārāntyaṁ)}]{\ml{വാരാന്ത്യം} (vārāntyaṁ) — the weekend}

\label{lesson:ML-C195-varantyam}



\begin{quote}
Before the new one: say the Malayalam for every day, then the Malayalam for a date.
\end{quote}

\subsection*{You'll want to know: \ml{വാരാന്ത്യം}}
\textbf{\ml{വാരാന്ത്യം}} — \emph{vārāntyaṁ} — \textquotedblleft{}the weekend\textquotedblright{}.

\begin{cousinweb}
\textbf{\ml{വാരാന്ത്യം}} (\emph{vārāntyaṁ}) is built from Sanskrit parts: \emph{vāra}, \textquotedblleft{}day of the week\textquotedblright{} (CDIAL 11547), and \emph{anta}, \textquotedblleft{}end\textquotedblright{} (CDIAL 347): the week's end.
\end{cousinweb}

\begin{grammarlens}[title={What you've built}]
One more word for this chapter.
\end{grammarlens}

\subsection*{Your turn}
\begin{itemize}
\raggedright
  \item \emph{Say it:} \emph{vārāntyaṁ}
  \item \emph{Say it:} \emph{vārāntyaṁ}, once more
  \item \emph{Say it:} say \emph{vārāntyaṁ}
  \item \emph{Recall:} say the Malayalam for every day, then the Malayalam for a date, then say \emph{vārāntyaṁ} again
\end{itemize}

\begin{tcolorbox}[breakable,colback=teal!4,colframe=teal!35!black,title={Before you move on}]
What does \ml{വാരാന്ത്യം} mean? (\textquotedblleft{}The weekend\textquotedblright{}.) Say it once more.
\end{tcolorbox}
\section[\texorpdfstring{\ml{പെട്ടെന്ന്} (peṭṭennŭ)}{പെട്ടെന്ന് (peṭṭennŭ)}]{\ml{പെട്ടെന്ന്} (peṭṭennŭ) — quickly, suddenly}

\label{lesson:ML-C195-pettennu}



\begin{quote}
Before the new one: say the Malayalam for a date, then the Malayalam for the weekend.
\end{quote}

\subsection*{You'll want to know: \ml{പെട്ടെന്ന്}}
\textbf{\ml{പെട്ടെന്ന്}} — \emph{peṭṭennŭ} — \textquotedblleft{}quickly, suddenly\textquotedblright{}.

\begin{cousinweb}
\textbf{\ml{പെട്ടെന്ന്}} (\emph{peṭṭennŭ}) is an expressive sound for suddenness followed by \emph{ennŭ}, \textquotedblleft{}saying\textquotedblright{} (DEDR 868): \textquotedblleft{}with a \emph{peṭṭ}\textquotedblright{}.
\end{cousinweb}

\begin{grammarlens}[title={What you've built}]
One more word for this chapter.
\end{grammarlens}

\subsection*{Your turn}
\begin{itemize}
\raggedright
  \item \emph{Say it:} \emph{peṭṭennŭ}
  \item \emph{Say it:} \emph{peṭṭennŭ}, once more
  \item \emph{Say it:} say \emph{peṭṭennŭ}
  \item \emph{Recall:} say the Malayalam for a date, then the Malayalam for the weekend, then say \emph{peṭṭennŭ} again
\end{itemize}

\begin{tcolorbox}[breakable,colback=teal!4,colframe=teal!35!black,title={Before you move on}]
What does \ml{പെട്ടെന്ന്} mean? (\textquotedblleft{}Quickly, suddenly\textquotedblright{}.) Say it once more.
\end{tcolorbox}
\section[\texorpdfstring{\ml{മുൻകൂട്ടി} (munkūṭṭi)}{മുൻകൂട്ടി (munkūṭṭi)}]{\ml{മുൻകൂട്ടി} (munkūṭṭi) — in advance}

\label{lesson:ML-C195-munkutti}



\begin{quote}
Before the new one: say the Malayalam for the weekend, then the Malayalam for quickly.
\end{quote}

\subsection*{You'll want to know: \ml{മുൻകൂട്ടി}}
\textbf{\ml{മുൻകൂട്ടി}} — \emph{munkūṭṭi} — \textquotedblleft{}in advance\textquotedblright{}.

\begin{cousinweb}
\textbf{\ml{മുൻകൂട്ടി}} (\emph{munkūṭṭi}) is native Dravidian, \textquotedblleft{}put together beforehand\textquotedblright{}: \emph{mun}, \textquotedblleft{}before\textquotedblright{} (DEDR 5020), and \emph{kūṭṭi}, from the verb \textquotedblleft{}to bring together\textquotedblright{} (DEDR 1882, with Tamil \emph{kūṭu}).
\end{cousinweb}

\begin{grammarlens}[title={What you've built}]
One more word for this chapter.
\end{grammarlens}

\subsection*{Your turn}
\begin{itemize}
\raggedright
  \item \emph{Say it:} \emph{munkūṭṭi}
  \item \emph{Say it:} \emph{munkūṭṭi}, once more
  \item \emph{Say it:} say \emph{munkūṭṭi}
  \item \emph{Recall:} say the Malayalam for the weekend, then the Malayalam for quickly, then say \emph{munkūṭṭi} again
\end{itemize}

\begin{tcolorbox}[breakable,colback=teal!4,colframe=teal!35!black,title={Before you move on}]
What does \ml{മുൻകൂട്ടി} mean? (\textquotedblleft{}In advance\textquotedblright{}.) Say it once more.
\end{tcolorbox}
